% !TeX root = Thesis.tex

\begin{theorem} \label{theorem::IP_gives_root}
  Let $G$ be a graph and consider the integer program whose variables $x_{ij} \in \{0,1\}$ encode a graph $H$ with adjacency matrix $A_H = (x_{ij})$. Then $G$ has a square root $H$ if and only if the integer  program has a feasible solution $x$. The feasible solution $x$ defines entries of the adjacency matrix of a graph $H$ that satisfies $H^2 = G$.
\end{theorem}
\begin{proof}
  Suppose $G$ has a square root $H$ and let $A_H$ be the adjacency matrix of $H$ and let  $x_{ij} = (A_H)_{ij}$. Construct the variables $x_{ij}, y_{ikj}, z_{ij}$ and the constant $g_{ij}$ as follows.
  \begin{enumerate}[label=\roman*)]
    \item Let $x_{ij} = 1$ if $ij \in E_H$ and $x_{ij}=0$ otherwise. 
    \item Let $y_{ikj} = x_{ik}x_{kj}$ for all $i,k,j$. 
    \item Let $\displaystyle z_{ij} = \sum_{k = 1}^{n} y_{ikj} + x_{ij} $ for all $i,j$. 
    \item Let $\displaystyle z_{ii} = \sum_{k  = 1}^{n} y_{iki} - \sum_{k = 1}^{n} x_{ik} =  \sum_{k = 1}^{n} x_{ik}^2 -  \sum_{k = 1}^{n} x_{ik} = 0 $
    since $x_{ij} \in \{0 , 1\} \implies x_{ik}^2 = x_{ik}$. 
    \item Let $g_{ij} = 1$ if $ij \in E_G$ and $g_{ij}=0$ otherwise. 
  \end{enumerate}
  Verifying constraints of the integer program for the variables 
  \begin{enumerate}[label=\roman*)]
    \item As $x_{ij} \in \{0, 1 \}$, then $y_{ikj}  \in \{0, 1 \}$, Equations~\eqref{eq::IP_x_variable}-\eqref{eq::IP_z_variable}
    and Equations~\eqref{eq::IP_x_upper_bound}-\eqref{eq::IP_y_upper_bound} are satisfied.  
    \item As $H$ is a simple directed graph, its adjacency matrix has zero diagonal and is symmetric. Therefore Equations~\eqref{eq::IP_symmetry_x}
    -\eqref{eq::IP_diagonal_x} are satisfied.  
    \item As $y_{ikj} = x_{ik}x_{kj}$ and all variables are binary. Therefore Equations~\eqref{eq::IP_product_upper_first}-
    \eqref{eq::IP_product_lower} are satisfied. 
    \item Equation~\eqref{eq::IP_indicator_arg} by definition is satisfied. 
    \item As $\displaystyle z_{ii} = \sum_{k = 1}^{n} x_{ik}x_{ki} - \sum_{k = 1}^{n} x_{ik}$, Equation~\eqref{eq::IP_fund_matrix} is satisfied. 
    \item As $\displaystyle z_{ij} \geq 1 \text{ if and only if } x_{ij} = 1 \text{ or there exists } k \text{ such that } x_{ik}=x_{kj}= 1$ and as $H^2 = G$, have $g_{ij} = 1 \text{ if and only if } 1 \leq d_H(i,j) \leq 2$. Thus $g_{ij} = 1 \text{ if and only if } z_{ij} \geq 1$. Therefore indicator constraints Equations~\eqref{eq::IP_indicator_upper}-\eqref{eq::IP_indicator_lower} hold. 
  \end{enumerate}
  All constraints of the integer program are satisfied by the constructed variables. Hence, the integer program is feasible.

  Suppose the integer program on graph $G$ has a feasible solution $x$, and let $x_{ij}$ define a graph $H$ by 
  \[
  ij \in E_H \text{ if and only if } x_{ij} = 1. 
  \]
  As $x_{ij}$ satisfies the constraints of the integer program, 
  \begin{enumerate}[label=\roman*)]
    \item Equation~\eqref{eq::IP_x_variable} and Equation~\eqref{eq::IP_x_upper_bound} gives $x_{ij} \in \{0 , 1\}$ for all $i,j$,
    \item Equation~\eqref{eq::IP_symmetry_x} and Equation~\eqref{eq::IP_diagonal_x} give $x_{ij} = x_{ji}$ and $x_{ii} = 0$,
  \end{enumerate}
  such that $x_{ij}$ defines a zero diagonal, symmetric matrix with all elements being either $0$ or $1$ which is a valid adjacency matrix of a graph. 
  Examining the auxiliary variables $y_{ikj}$ and $z_{ij}$, 
  \begin{enumerate}[label=\roman*)]
    \item Equations~\eqref{eq::IP_product_upper_first}-\eqref{eq::IP_product_lower} give the strict equality $y_{ikj} = x_{ik}x_{kj}$ for all $i,j,k$,
    \item Equation~\eqref{eq::IP_indicator_arg} gives $\displaystyle z_{ij} = \sum_{k = 1}^{n}y_{ikj} + x_{ij} = \sum_{k = 1}^{n}x_{ik}x_{kj} + x_{ij}$ for all $i \neq j$,
    \item Equation~\eqref{eq::IP_fund_matrix} gives $\displaystyle z_{ii} = \sum_{k = 1}^{n}y_{iki} + \sum_{n}^{k = 1}x_{ik} = \sum_{k = 1}^{n} x_{ik}x_{ki} - \sum_{n}^{k = 1}x_{ik} = 0$ for all $i$ since
    \[
    \sum_{k = 1}^{n} x_{ik}x_{ki} = \sum_{k = 1}^{n} x_{ik}^2 = \sum_{k = 1}^{n} x_{ik},
    \]
    by symmetry and since $x_{ij} \in \{0 , 1\}$. 
  \end{enumerate} 
  The indicator constraints Equations~\eqref{eq::IP_indicator_upper}-\eqref{eq::IP_indicator_lower} thus give for $i \neq j$
  \[
  g_{ij} = 1 \text{ if and only if } \sum_{k=1}^{n} x_{ik}x_{kj} + x_{ij} \geq 1  \text{ and }  g_{ij} = 0 \text{ if and only if } \sum_{k=1}^{n} x_{ik}x_{kj} + x_{ij} = 0. 
  \]
  Since $x_{ij} \in \{0,1\}$, this is equivalent to saying that $
  g_{ij} = 1$ if and only if $\Big(x_{ij} = 1 \Big) \lor \Big(\text{ there exists } k $ \newline
  $\text{ such that } x_{ik}= 1 \land x_{kj}=1\Big)$ and  $g_{ij} = 0$ if and only if $x_{ij} = 0$. 
  By definition of $H$ this means that $g_{ij} = 1 \text{ if and only if } \Big( ij \in E_H \Big) \lor \Big( \text{ there exists } k \text{ such that } ik,kj \in E_H\Big)$ which is equivalent to 
  \[
  g_{ij} = 1 \text{ if and only if } 1 \leq d_H(i,j) \leq 2. 
  \]
  Applying Lemma~\ref{lemma::G_equal_to_H^2} thus gives $ij \in E_G \text{ if and only if } ij \in E_H^2$. Therefore $G=H^2$. 
\end{proof}
